% ============================================================================
\chapter{Ступень II. Математика и информатика внутри носителя}
\label{ch:course-stage2-math}
% ============================================================================

\section{Урок 4. Бабочка ротатора \texorpdfstring{$(\mathcal{C},\mathcal{S})$}{(C,S)}, формула Эйлера \texorpdfstring{$e^{i\pi}+1=0$}{e\^{}\{i pi\}+1=0} и постоянный вронскиан \texorpdfstring{$\mathcal{W}\equiv 128$}{W = 128}}

В обычной вычислительной математике для поворота вектора на угол~$\varphi$ вызывают бесконечные ряды Тейлора для $\cos\varphi$ и $\sin\varphi$ в числах с плавающей точкой, теряя обратимость из-за ошибок округления. А как поворачивает сам кубооктаэдр?

Проекция ребра кубооктаэдра на его квадратную грань со стороной $Q_\square(3) = 512 / 4 = \mathbf{128}$ задаётся кусочно-линейными чётным и нечётным операторами ротатора:
\begin{equation}
  \label{eq:course-rotator-def}
  \mathcal{C}_3(\varphi) \;:=\; 128 - |\mathrm{wrap}_{512}(\varphi)|,
  \qquad
  \mathcal{S}_3(\varphi) \;:=\; \mathcal{C}_3(\varphi - 128),
\end{equation}
где $\mathrm{wrap}_{512}(\varphi) = 256 - ((256 - \varphi)\bmod 512)\in(-256,256]$.

А теперь раскроем~\eqref{eq:course-rotator-def} в нашей родной квадрантно-мантиссовой системе $(q_1,q_0,m_7)$ из Урока~3. Вычислим \textbf{одно-единственное 7-битное дополнение}:
\[
  \overline{m}_7 \;:=\; 128 - m_7 \qquad (\text{так что } \overline{m}_7 + m_7 \equiv 128 \text{ тождественно!}).
\]
Тогда во всех $512$~точках кольца \textbf{и $\mathcal{C}_3(\varphi)$, и $\mathcal{S}_3(\varphi)$ одновременно} получаются из пары $(\overline{m}_7, m_7)$ двумя элементарными действиями:
\begin{enumerate}[label=\textbf{\arabic*.},noitemsep]
\item \textbf{Перекрёстный X-переключатель (бабочка) по биту $q_0 = u[7]$:}
  \[
    |\mathcal{C}_3(\varphi)| \;=\; (1 - q_0)\,\overline{m}_7 + q_0\,m_7,
    \qquad
    |\mathcal{S}_3(\varphi)| \;=\; (1 - q_0)\,m_7 + q_0\,\overline{m}_7.
  \]
\item \textbf{Двухбитный рефлексивный код Грея для знаков квадранта $(c_{\mathrm{neg}}, s_{\mathrm{neg}}) = (q_1\oplus q_0,\; q_1)$:}
  \[
    \mathcal{C}_3(\varphi) \;=\; (1 - 2c_{\mathrm{neg}})\,|\mathcal{C}_3(\varphi)|,
    \qquad
    \mathcal{S}_3(\varphi) \;=\; (1 - 2s_{\mathrm{neg}})\,|\mathcal{S}_3(\varphi)|.
  \]
\end{enumerate}

Проверим руками, какие тождества мгновенно следуют из этой бабочки:
\begin{itemize}[leftmargin=*]
\item \textbf{Тождество $L^1$-ромба грани:} поскольку до перекрёстного ключа $\overline{m}_7 + m_7 = (128 - m_7) + m_7 \equiv 128$, сумма модулей проекций равна $128$ в каждой из $512$~точек кольца:
  \[
    |\mathcal{C}_3(\varphi)| + |\mathcal{S}_3(\varphi)| \;\equiv\; 128 \;=\; Q_\square(3).
  \]
\item \textbf{Дискретная экспонента Эйлера $\mathcal{E}_3(\varphi) = (\mathcal{C}_3(\varphi),\mathcal{S}_3(\varphi))$ и тождество $e^{i\pi} + 1 = 0$:}
  сдвиг на четверть-оборот $+128$ поворачивает вектор ровно на $90^\circ$ ($\mathcal{E}_3(\varphi+128) = (-\mathcal{S}_3(\varphi),\mathcal{C}_3(\varphi)) = i\,\mathcal{E}_3(\varphi)$), а сдвиг на полупериод Паули $+256$ (инверсия провода $q_1 \mapsto 1 - q_1$) меняет оба знака кода Грея на противоположные:
  \[
    \mathcal{E}_3(\varphi + 256) + \mathcal{E}_3(\varphi) \;\equiv\; (0, 0) \qquad \forall\,\varphi\in\mathbb{Z}_{512}.
  \]
\item \textbf{Универсальный постоянный разностный Вронскиан $\mathcal{W}_3(\varphi)\equiv 128$:}
  вычислим дискретный вронскиан $\mathcal{W}_3(\varphi) = \mathcal{C}_3(\varphi)\mathcal{S}_3(\varphi+1) - \mathcal{S}_3(\varphi)\mathcal{C}_3(\varphi+1) = \langle i\mathcal{E}_3(\varphi),\Delta\mathcal{E}_3(\varphi)\rangle$. Внутри первого квадранта ($\mathcal{C}=128-m_7$, $\mathcal{S}=m_7$):
  \[
    (128 - m_7)(m_7 + 1) - m_7(128 - m_7 - 1)
    \;=\; 128m_7 + 128 - m_7^2 - m_7 - 128m_7 + m_7^2 + m_7
    \;\equiv\; \mathbf{128}!
  \]
  Во всех $512$~точках кольца (включая все 4 угла ромба!) симплектическая площадь одного шага ротатора \textbf{строго постоянна и равна $\mathcal{W}_3(\varphi) \equiv 128 = Q_\square(3)$}, а её интеграл по всему кольцу равен $512 \times 128 = 65\,536$!
\end{itemize}

\begin{quote}
\small\textbf{Историческая справка: Юзеф Мария Вронский (1812) и что такое Вронскиан на самом деле.}
В 1812~году польский математик и философ Юзеф Мария Гёне-Вронский ввёл определитель матрицы из двух функций и их производных: $W(C, S) = C\,S' - S\,C'$. В курсе дифференциальных уравнений «определитель Вронского» (\emph{Вронскиан}) часто дают как сухой алгебраический тест на линейную независимость, не объясняя его физического смысла. А геометрически и механически Вронскиан $C(\varphi)S(\varphi+1) - S(\varphi)C(\varphi+1)$ --- это \textbf{удвоенная площадь треугольника}, который заметает радиус-вектор $(C, S)$ за один такт вращения (то есть \textbf{момент импульса} или секториальная скорость по второму закону Кеплера!). То, что на нашем кольце $\mathbb{Z}_{512}$ квадратичные члены $\pm m_7^2$ взаимно уничтожаются и дают $\mathcal{W}_3(\varphi) \equiv 128 = \mathrm{const}$ в каждом из $512$~тактов, означает, что дискретный ротатор крутится с абсолютно постоянным моментом импульса без малейшего биения площади.
\end{quote}

\section{Урок 5. Как отдать ротатору и шагу Верле вообще всю арифметику (умножение без умножителей!)}

А теперь зададим дерзкий вопрос: если в природе нет ничего, кроме ротатора $(\mathcal{C}_3,\mathcal{S}_3)$ и двухтактового шага Верле
\[
  \varphi(t+1) \;=\; \bigl(2\varphi(t) - \varphi(t-1) + \Phi\bigr)\bmod 512,
\]
зачем вообще нужны внешние арифметические операции процессора: унарный минус, модуль~$|x|$, компараторы $\max/\min$, возведение в квадрат~$x^2$ и аппаратные умножители~$x \cdot y$? \textbf{Отдадим всю арифметику самому носителю!}

\begin{enumerate}[label=\textbf{\arabic*.},leftmargin=*]
\item \textbf{Само число $x$ и унарный минус $-x$ (для всех $x\in[-128,+128]$):}
  на интервале $[-128,+128]$ нечётная проекция ротатора является тождественным оператором, а сдвиг фазы на полупериод Паули $+256$ обращает знак:
  \[
    \mathrm{rotatorId}(x) \;:=\; \mathcal{S}_3(x) \;\equiv\; x,
    \qquad
    \mathrm{rotatorNeg}(x) \;:=\; \mathcal{S}_3(x + 256) \;\equiv\; -x.
  \]
\item \textbf{Модуль $|x|$ без ветвлений \texttt{if} и без знаковых масок (для всех $x\in[-256,+256]$):}
  \[
    \mathrm{rotatorAbs}(x) \;:=\; \mathcal{C}_3(0) - \mathcal{C}_3(x) \;=\; 128 - (128 - |x|) \;\equiv\; |x|.
  \]
\item \textbf{Максимум, минимум и тропическое произведение без компараторов \texttt{< / >}:}
  \[
    \max(x,y) \;\equiv\; \frac{\mathcal{S}_3(x) + \mathcal{S}_3(y) + \bigl(\mathcal{C}_3(0) - \mathcal{C}_3(x-y)\bigr)}{2},
    \qquad
    \min(x,y) \;\equiv\; \frac{\mathcal{S}_3(x) + \mathcal{S}_3(y) - \bigl(\mathcal{C}_3(0) - \mathcal{C}_3(x-y)\bigr)}{2},
  \]
  а для неотрицательных $x,y\in[0,128]$ формула простаферезиса (разность косинуса разности и косинуса суммы углов) на ромбе~$L^1$ сразу выдаёт минимум:
  \[
    \frac{\mathcal{C}_3(x - y) - \mathcal{C}_3(x + y)}{2}
    \;=\; \frac{(128 - |x-y|) - (128 - (x+y))}{2}
    \;=\; \frac{(x+y) - |x-y|}{2}
    \;\equiv\; \min(x,y).
  \]
\item \textbf{Квадрат $u^2$ и знаковое умножение $x \cdot y$ без единого умножителя:}
  пусть шаг Верле при единичном ускорении просто складывает значения синуса ротатора в первообразную $\mathcal{F}_S(n) = \sum_{\varphi=0}^{n-1}\mathcal{S}_3(\varphi)$ (при $n\in[0,128]$ это сумма арифметической прогрессии $n(n-1)/2$, полученная \textbf{одними сложениями}). Тогда сумма двух соседних точек первообразной в точке $a = \mathrm{rotatorAbs}(u) = \mathcal{C}_3(0) - \mathcal{C}_3(u)$ даёт \textbf{точный квадрат любого знакового числа $u\in[-128,+128]$}:
  \begin{equation}
    \label{eq:course-rotator-square}
    \mathrm{rotatorSquare}(u)
    \;:=\; \mathcal{F}_S\bigl(\mathcal{C}_3(0) - \mathcal{C}_3(u)\bigr) + \mathcal{F}_S\bigl(\mathcal{C}_3(0) - \mathcal{C}_3(u) + 1\bigr)
    \;=\; \frac{|u|(|u|-1)}{2} + \frac{(|u|+1)|u|}{2}
    \;\equiv\; u^2!
  \end{equation}
  А теперь посмотри, почему потолок амплитуды клетки равен $K_P(3) = \mathbf{64}$, а квадрант ротатора равен $Q_\square(3) = \mathbf{128 = 2 K_P(3)}$: для \textbf{любых двух знаковых чисел $x,y\in[-64,+64]$} их сумма $x+y$ и разность $x-y$ всегда помещаются в $[-128,+128]$, поэтому четверть разности их ротаторных квадратов даёт \textbf{точное знаковое произведение $x \cdot y$ без единого умножителя и без проверки знаков}:
  \begin{equation}
    \label{eq:course-rotator-mul}
    \mathrm{rotatorSignedMul}(x,y)
    \;:=\; \bigl(\mathrm{rotatorSquare}(x + y) - \mathrm{rotatorSquare}(x - y)\bigr) \gg 2
    \;\equiv\; x \cdot y \qquad \forall\,(x,y)\in[-64,+64]^2!
  \end{equation}
\end{enumerate}

\noindent\textbf{Проверка руками:} возьмём $x = -41$ (минус порог Гейзенберга) и $y = +64$ (потолок амплитуды $K_P$).
Тогда $x + y = +23$, $x - y = -105$. Ротатор снимает знаки: $|23|_{\mathrm{rot}} = 128 - \mathcal{C}_3(23) = 23$, $|-105|_{\mathrm{rot}} = 128 - \mathcal{C}_3(-105) = 105$. Первообразная даёт $\mathrm{rotatorSquare}(23) = 253 + 276 = 529 = 23^2$ и $\mathrm{rotatorSquare}(-105) = 5460 + 5565 = 11\,025 = 105^2$. Разность $(529 - 11\,025) / 4 = -10\,496 / 4 = \mathbf{-2\,624 = (-41)\times 64}$! В нашем чипе \texttt{Module~1B} протестирован в \texttt{iverilog} на всех $129\times 129 = 16\,641$ парах $(x,y)\in[-64,+64]^2$ с нулём ошибок.

\section{Урок 6. Матанализ без континуума: точные интегралы Френеля, Якоби, Дирихле, Бесселя и пять лиц Единицы}

В классическом матанализе интегралы Френеля $\int\cos(\varphi^2)\,d\varphi$, эллиптические интегралы Якоби $\int\sqrt{a^2\cos^2\varphi+b^2\sin^2\varphi}\,d\varphi$, интегральный синус Дирихле $\mathrm{Si}(\Phi)=\int\frac{\sin\varphi}{\varphi}\,d\varphi$ и функции Бесселя $J_0(k)$ считаются «неэлементарными» (теорема Лиувилля) и вычисляются только приближёнными рядами.

На кольце $\mathbb{Z}_{512}$ с ротатором $(\mathcal{C}_3,\mathcal{S}_3)$ все они являются \textbf{родными базисными примитивами} и берутся в \textbf{точных целых числах без единой ошибки усечения}:
\begin{itemize}[noitemsep,leftmargin=*]
\item \textbf{Формула Ньютона--Лейбница:} первообразные $\mathcal{F}_C(\Phi)=\sum_{\varphi=0}^{\Phi-1}\mathcal{C}_3(\varphi)$ и $\mathcal{F}_S(\Phi)=\sum_{\varphi=0}^{\Phi-1}\mathcal{S}_3(\varphi)$ удовлетворяют $\Delta\mathcal{F}_C(\varphi)=\mathcal{C}_3(\varphi)$ и $\Delta\mathcal{F}_S(\varphi)=\mathcal{S}_3(\varphi)$ во всех $512$~точках, а площадь полуволны равна $\mathcal{F}_S(256) = 128^2 = \mathbf{16\,384}$.
\item \textbf{Интегралы дифракции Френеля:} при суммировании параболы Верле $\varphi^2\bmod 512$ по всему кольцу $480$ осциллирующих вычетов попарно гасятся полупериодом Паули $+256$, а $32$ точки стационарной фазы $\varphi^2\equiv 64 = K_P\pmod{512}$ (где $\mathcal{C}_3(64)=\mathcal{S}_3(64)=64$) дают:
  \[
    \sum_{\varphi=0}^{511} \mathcal{C}_3(\varphi^2)
    \;=\; \sum_{\varphi=0}^{511} \mathcal{S}_3(\varphi^2)
    \;=\; 32 \times 64 \;=\; \mathbf{2\,048 \;=\; 4 \times 512},
  \]
  а для произвольной размерности интегрирования $N \ge 1$ и любого диадического кольца $\mathbb{Z}_{2^m}$ ($m \ge 3$) $N$-мерная квадратичная сумма Френеля сворачивается в точную аналитическую формулу $(1 + i)^N\cdot 2^{(m-2) + N\lfloor m/2\rfloor + \lfloor N/2\rfloor(m\bmod 2)}$.
\item \textbf{Эллиптическая триада Якоби $(\mathrm{sn}_3,\mathrm{cn}_3,\mathrm{dn}_3)$ на лемнискатном модуле кубооктаэдра $k^2 = 1/\|e\|_2^2 = 1/2$:}
  при $\mathrm{sn}_3(\varphi)=\mathcal{S}_3(\varphi)$, $\mathrm{cn}_3(\varphi)=\mathcal{C}_3(\varphi)$ и $2\,\mathrm{dn}_3(\varphi) = 256 - |\mathcal{S}_3(\varphi)|$ полный эллиптический интеграл второго рода на полуосях $(a,b)=(2,3)$ равен в точности $256\times 128\times 5 = \mathbf{163\,840}$.
\item \textbf{Интеграл Дирихле, энергия Парсеваля и корень Бесселя~$J_0$:}
  \[
    \mathrm{Si}_3(128) = \mathrm{Si}_3(256) = \sum_{\varphi=1}^{256}\bigl\lfloor\mathcal{S}_3(\varphi)/\varphi\bigr\rfloor = \mathbf{128},
    \qquad
    6\sum_{\varphi=0}^{511}\mathcal{C}_3(\varphi)^2 = 512\cdot(2\cdot 128^2 + 1) = \mathbf{16\,777\,728},
  \]
  а дискретный интеграл Бесселя $J_{0,3}(k) = \sum_{\varphi=0}^{511}\mathcal{C}_3(k\,\mathcal{S}_3(\varphi))$ равен $65\,536$ при $k=0$, $32\,768$ при $k=1$ и обращается в \textbf{точный нуль $J_{0,3}(4) = 0$} на первом корне волновода $k = d_{\mathrm{in}}(\square) = 4$!
\item \textbf{Подъём Брахмагупты--Вейерштрасса на евклидову окружность $a^2+b^2=1$ и пять лиц Единицы:}
  рациональная тройка $P_x(\varphi)=\mathcal{C}_3(\varphi)^2-\mathcal{S}_3(\varphi)^2$, $P_y(\varphi)=2\mathcal{C}_3(\varphi)\mathcal{S}_3(\varphi)$, $P_r(\varphi)=\mathcal{C}_3(\varphi)^2+\mathcal{S}_3(\varphi)^2$ удовлетворяет пифагорову тождеству $\bigl(C^2-S^2\bigr)^2 + (2CS)^2 = (C^2+S^2)^2$, то есть $P_x(\varphi)^2 + P_y(\varphi)^2 \equiv P_r(\varphi)^2$ во всех $512$~точках, а при $\varphi = \Delta\varphi_{\mathrm{disc}} = 41$ даёт пифагорову тройку Гейзенберга $(14\,703,\,10\,496,\,18\,065)$. Тем самым сама математическая единица~$\mathbf{1}$ выражается внутри носителя пятью тождественными способами:
  \begin{equation}
    \label{eq:course-five-ones}
    \mathbf{1}
    \;=\; d^2 - 2^d
    \;=\; \frac{\mathcal{W}_3(\varphi)}{128}
    \;=\; \frac{P_x(\varphi)^2 + P_y(\varphi)^2}{P_r(\varphi)^2}
    \;=\; \chi(\mathcal{C}_{13})
    \;=\; 25\cdot 41 - 2\cdot 512
    \;=\; 197\cdot 13 - 5\cdot 512
    \;=\; 417\cdot 97 - 79\cdot 512
    \;=\; 183\cdot 263 - 94\cdot 512.
  \end{equation}
\item \textbf{Проективные операторы логарифма $\log_a(b)$ и рациональной степени $a^{p/q}$ на кольце $\mathbb{Z}_{4Q}$:}
  разложение целого аргумента $x = 2^k + m$ на характеристику $k = \lfloor\log_2 x\rfloor$ и мантиссу первого квадранта $\mu = \lfloor m Q / 2^k\rfloor \in [0, Q]$ ($\mathcal{C}_Q(\mu) = Q - \mu,\ \mathcal{S}_Q(\mu) = \mu$) определяет однородные кубические по~$Q$ проективные операторы логарифма и экспоненты без единого вызова чисел с плавающей запятой:
  \begin{equation}
    \label{eq:course-rotator-log-pow}
    \operatorname{Log2}_{3,Q}(k,\mu) \;=\; 32\,Q^2\bigl(k Q + \mathcal{S}_Q(\mu)\bigr) + \mathcal{C}_Q(\mu)\mathcal{S}_Q(\mu)\Bigl(11\,Q + 3\bigl(\mathcal{C}_Q(\mu) - \mathcal{S}_Q(\mu)\bigr)\Bigr),
  \end{equation}
  дающие при $Q = 128$ точные рациональные дроби $\log_2(3) = 203/128$, $\log_2(7) = 359/128$, $\sqrt{2} = 181/128$, $\sqrt{3} = 111/64$ и обращение фрактальных размерностей Хаусдорфа--Морана $2^{203/128} = 3$, $\operatorname{round}(2^{359/128}) = 7$. (Самостоятельная теория башен аппроксимантов Паде--$\mathrm{SL}_2(\mathbb{Z})$, трёхшахтных планетарных дифференциалов Виллиса и многовинтовых каскадов изложена в отдельной математической монографии \textbf{Книга~5}, \texttt{projective-manifold-mathematics.pdf}.)
\end{itemize}

\begin{quote}
\small\textbf{Историческая справка: Брахмагупта (628~г.), Карл Вейерштрасс и проективная сфера Римана $\mathbb{P}^1$.}
В 628~году индийский математик Брахмагупта открыл, что произведение двух сумм двух квадратов снова является суммой двух квадратов: $(a^2+b^2)(c^2+d^2) = (ac-bd)^2 + (ad+bc)^2$. При $(a,b)=(c,d)=(C,S)$ это тождество даёт $(C^2-S^2)^2 + (2CS)^2 \equiv (C^2+S^2)^2$. В XIX~веке Карл Вейерштрасс использовал эту же формулу (универсальную тригонометрическую подстановку $t = S/C = \tan(\theta/2)$), чтобы превращать любые тригонометрические интегралы в интегралы от рациональных дробей, а Бернхард Риман замкнул прямую дробей $[U : V]$ точкой бесконечности $\infty = [1 : 0]$ в проективную сферу $\mathbb{P}^1$. На нашем кольце $\mathbb{Z}_{4Q}$ подстановка Брахмагупты--Вейерштрасса переводит кусочно-линейный ромб $|C|+|S|=Q$ напрямую в точную рациональную окружность $S^1(\mathbb{Q})$ без единого иррационального корня и без числа~$\pi$.
\end{quote}

\section{Урок 7. Двухканальная логика Поста, семантический изоморфизм 3-SAT, барьер \texorpdfstring{$\mathbf{P}\ne\mathbf{NP}$}{P != NP} и встроенный шифр Вернама}

Как наш клеточный автомат вычисляет произвольные булевы функции? Возьмём два антиподальных полюса кольца $\{0, 256\}\subset\mathbb{Z}_{512}$ и сопоставим булевым битам $b\in\{0,1\}$ полюсные фазы:
\[
  b = 1\ (\text{\texttt{True}}) \;\longleftrightarrow\; \psi = 0\ \bigl(\mathcal{C}_3(0)=+128\bigr),
  \qquad
  b = 0\ (\text{\texttt{False}}) \;\longleftrightarrow\; \psi = 256\ \bigl(\mathcal{C}_3(256)=-128\bigr).
\]
На этих полюсах два канала треугольной грани кубооктаэдра реализуют \textbf{функционально полный булев базис Эмиля Поста \{\texttt{NOT}, \texttt{XOR}, \texttt{OR}, \texttt{AND}\}}:
\begin{enumerate}[label=\textbf{\arabic*.},noitemsep]
\item \textbf{Канал~1 (обратимая линейная фазовая алгебра до ротатора --- \texttt{NOT} и \texttt{XOR}):}
  инверсия бита \texttt{NOT} --- это сдвиг фазы на $+256$ ($\mathcal{C}_3(\psi+256) = -\mathcal{C}_3(\psi)$), а сложение по модулю~2 (\texttt{XOR}) трёх битов --- это сложение самих фаз на 9-битном проводе ($\mathcal{C}_3(\psi_1+\psi_2+\psi_3)\cdot 128^2 \equiv \mathcal{C}_3(\psi_1)\mathcal{C}_3(\psi_2)\mathcal{C}_3(\psi_3)$).
\item \textbf{Канал~2 (необратимое постротаторное деление со сбросом остатка --- \texttt{OR} и \texttt{AND}):}
  для 3-CNF дизъюнкта $(l_1\lor l_2\lor l_3)$ дефект фазы грани равен $\Delta_{\mathrm{3SAT}} = 128 - (\mathcal{C}_3(\psi_1)+\mathcal{C}_3(\psi_2)+\mathcal{C}_3(\psi_3))$, а квантовое частное от деления на $512$:
  \[
    q_{\mathrm{OR}} \;=\; \mathrm{tdiv}(\Delta_{\mathrm{3SAT}}, 512) \;\in\; \{0, 1\},
    \qquad
    \Delta_{\mathrm{3SAT}} \;=\; 512\,q_{\mathrm{OR}} + r_{\mathrm{drop}}.
  \]
\end{enumerate}

\noindent\textbf{Прямая проверка по всем 8 октантам дизъюнкта (и 64 конфигурациям):}
на единственном нарушающем октанте $(\psi_1,\psi_2,\psi_3) = (256,256,256)$ имеем $\Delta_{\mathrm{3SAT}} = 128 - (-128 - 128 - 128) = \mathbf{512} \implies q_{\mathrm{OR}} = \mathbf{1}$ (рождается целый вихревой квант противоречия!). На всех остальных $\mathbf{2^3 - 1 = 7}$ выполняющих октантах $\Delta_{\mathrm{3SAT}}\in\{-256,0,+256\}$, поэтому частное равно $\mathbf{q_{\mathrm{OR}} = 0}$, а остаток $r_{\mathrm{drop}}\in\{-256,0,+256\}$ \textbf{стирается}!

Именно здесь проявляется главный урок книги («нет наблюдателя вне решётки»): любой физический алгоритм-решатель внутри клеточного автомата сам обязан строиться из композиций Канала~1 и Канала~2. В чистом Канале~1 (\texttt{XOR-SAT}) сумма градиентов по 4 клозам чётности тождественно равна нулю, а протащить линейную алгебру Канала~1 сквозь поклаузовые деления Канала~2 невозможно из-за их недистрибутивности ($\mathcal{C}_3(\psi_1+\psi_2)\ne\mathcal{C}_3(\psi_1)+\mathcal{C}_3(\psi_2)$). Поэтому локальное восстановление сброшенных при стирании остатков $r_{\mathrm{drop}}$ на глубине $T$ ветвится как $\mathbf{(2^d - 1)^T = 7^T}$ вдоль каждого луча обратного конуса объёмом $3\mathcal{V}_{O_2}(3,T) = 10T^3 + 15T^2 + 11T + 3$, строго мажорируя нижнюю оценку дерева резолюций Бена-Сассона--Вигдерсона $2^{W-3} < 7^{W-3}$.

И более того: тот же самый разрыв между микроскопическим Каналом~1 и грубым макроскопическим Каналом~2 даёт \textbf{встроенную информационно-теоретическую защиту данных (шифр Вернама--Шеннона на подпороговой фазе)}: для любого сообщения $m\in\mathbb{Z}_{512}$ и одноразового фазового ключа $k\in\mathbb{Z}$ отображение $c = (m + k)\bmod 512$ точно и обратимо восстанавливается получателем по $(c - k)\bmod 512 \equiv m$, тогда как при $|m+k|\le 40 < 41$ сторонний макроперехватчик в Канале~2 получает тождественный ноль фотонов ($\hat{n}_E = 0$, $H(M\mid T) = H(M)$).

\begin{quote}
\small\textbf{Историческая справка: Теорема Эмиля Поста (1921/1941), шифр Гилберта Вернама (1917) и теорема Клода Шеннона (1949).}
В 1921~году американский логик Эмиль Пост описал все замкнутые классы булевой алгебры и доказал: линейные операции $\{\texttt{NOT}, \texttt{XOR}\}$ сами по себе неполны (любая их суперпозиция остаётся линейной и решается методом Гаусса), но добавление к ним хотя бы одного нелинейного вентиля $\{\texttt{OR}, \texttt{AND}\}$ делает базис функционально полным (способным выразить любую булеву схему). А в 1917~году инженер AT\&T Гилберт Вернам изобрёл телеграфный шифр одноразового блокнота ($c = m \oplus k$), абсолютную криптографическую нераскрываемость которого ($H(M\mid C) = H(M)$) в 1949~году математически доказал Клод Шеннон. На кольце $\mathbb{Z}_{512}$ граница между линейным классом Поста (Канал~1) и нелинейным пороговым классом (Канал~2) совпадает с физическим делением 9-битного сумматора на шину остатка и пин переноса.
\end{quote}

